%% kapa: exam=1 round=33 year=2022 period=1 form=A subject=civil_law subject_ko=민법 seq=24 number=24
%% source: Q-Net 감정평가사 기출문제 (cst00309) artlSeq=5209997 file=raw/1008/5209997/2022년도 제33회 감정평가사 1차시험 1교시(민법,경제학,부동산학).pdf
%% conversion: kordoc 4.24.51 -> tools/gen1.py (auto); review=proofread-sonnet; flags=none
\begin{question}{24}
첨부에 관한 설명으로 옳지 않은 것은? (다툼이 있으면 판례에 따름)
\begin{choices}
  \item 주종을 구별할 수 있는 동산들이 부합하여 분리에 과다한 비용을 요할 경우, 그 합성물의 소유권은 주된 동산의 소유자에게 속한다.
  \item 타인이 권원에 의하여 부동산에 부속시킨 동산이 그 부동산과 분리되면 경제적 가치가 없는 경우, 그 동산의 소유권은 부동산 소유자에게 속한다.
  \item 양도담보권의 목적인 주된 동산에 甲소유의 동산이 부합되어 甲이 그 소유권을 상실하는 손해를 입은 경우, 특별한 사정이 없는 한 甲은 양도담보권자를 상대로 보상을 청구할 수 있다.
  \item 타인의 동산에 가공한 경우, 가공으로 인한 가액의 증가가 원재료의 가액보다 현저히 다액인 때에는 가공자의 소유로 한다.
  \item 건물의 증축 부분이 기존 건물에 부합하여 기존 건물과 분리해서는 별개의 독립물로서의 효용을 갖지 못하는 경우, 기존 건물에 대한 경매절차에서 경매목적물로 평가되지 않았더라도 매수인은 부합된 증축 부분의 소유권을 취득한다.
\end{choices}
\end{question}
\answer{3}
